%%%%%%%%%%%%%%%%%
% This is the resume editor's preamble for rt-awesome-cv.cls, its modified version of
% awesome-cv.cls, the class file of Awesome CV by Claud D. Park (posquit0.bj@gmail.com), which
% is at https://github.com/posquit0/Awesome-CV. The class is under the LaTeX Project Public
% License 1.3c, and rt-awesome-cv.cls says what the editor changed in it. This preamble is the
% editor's own. It takes from the example resume.tex that comes with the class, which is under
% CC BY-SA 4.0, only the class options, the margins and the settings that restate the class's
% defaults, and none of its text.
%%%%%%%%%%%%%%%%

% This is the editor's port of Awesome CV. The fonts, the colours and the layout are the class's,
% as its example sets them, and the preamble changes the example only as follows. The font
% encoding is OT1 until the class has loaded, and TU after it, so that the format the editor
% dumps from this preamble holds no native font; XeTeX loads each font the first time the
% document selects it. The margins control sets the margins after the format, the same on every
% side, and its default is half an inch, which of what the control offers comes nearest the
% example's margins taken together, 1.4 cm at the sides, 0.8 cm at the top and 1.8 cm at the
% bottom, so the example's geometry below sets only the space for the footer. The accent colour
% is named, rtaccent, and the example's red, awesome-red, is its default; the accent control
% defines it again and sets the class's awesome from it. The class sets its text in sizes of its
% own, fixed in points, from its 11 points, while the font size control takes half points from 9
% to 12 as an ordinary compile: \rtfontsize sets every size of the class's and of article's
% size11.clo scaled by the size asked for over 11 points, with the space between paragraphs that
% goes with it, and \rtdensity scales the space between sections, entries and bullets. At 11
% points and normal density, the controls' defaults, every size and every space is the class's.
% The font control sets the text in Source Sans Pro, Roboto or Lato, and the header stays
% Roboto. A list section sets its items as the class sets an entry's, in a paragraph of its own.
% And the file ends where the body would begin, because the editor generates the body from the
% resume.

% OT1 until the class has loaded, as decision 0003 has a XeTeX port, so that nothing selects a
% native font before the format is dumped; the document's encoding is TU, XeTeX's, from below.
\RequirePackage[OT1]{fontenc}
\documentclass[11pt,a4paper]{rt-awesome-cv}

% The example's page. The margins control sets the margins after the format.
\geometry{left=1.4cm, top=.8cm, right=1.4cm, bottom=1.8cm, footskip=.5cm}

% The fonts. The document's text is TU, in Source Sans Pro, as sourcesanspro's default option
% makes it. \rtbodyfonts{<family>}{<light family>} is the font control's: the text in another
% of the families the package declares, the header staying Roboto.
\renewcommand{\encodingdefault}{TU}
\renewcommand{\familydefault}{SourceSans}
\newcommand*\rtbodyfonts[2]{%
  \renewcommand*{\bodyfont}{\fontencoding{TU}\fontfamily{#1}\selectfont}%
  \renewcommand*{\bodyfontlight}{\fontencoding{TU}\fontfamily{#2}\selectfont}%
  \renewcommand*{\familydefault}{#1}}

% The font size control's sizes, and the density control's spacing, as the note at the top says.
% \rt@size is the size in scaled points, 11 points until \rtfontsize sets another, and
% \rt@pt{<points>} is a length scaled by it over 11 points and rounded to a hundredth of a point,
% TeX arithmetic throughout. A length to two decimals rounds back to what TeX reads for it, so at
% 11 points every length is the class's.
\makeatletter
\newcommand*\rt@size{720896}
\newcommand*\rt@pt[1]{%
  \the\dimexpr\numexpr\dimexpr#1\p@*\rt@size/720896\relax*100/65536\relax\p@/100\relax}
% size11.clo's sizes, with the display skips it sets with them, each length through \rt@pt. The
% lists' spacing is parskip's, which the class loads: no space above or between items, and the
% space between paragraphs within one.
\DeclareRobustCommand\normalsize{%
  \@setfontsize\normalsize{\rt@pt{10.95}}{\rt@pt{13.6}}%
  \abovedisplayskip \rt@pt{11}\@plus\rt@pt{3}\@minus\rt@pt{6}%
  \abovedisplayshortskip \z@ \@plus\rt@pt{3}%
  \belowdisplayshortskip \rt@pt{6.5}\@plus\rt@pt{3.5}\@minus\rt@pt{3}%
  \belowdisplayskip \abovedisplayskip
  \let\@listi\@listI}
\DeclareRobustCommand\small{\@setfontsize\small{\rt@pt{10}}{\rt@pt{12}}}
\DeclareRobustCommand\footnotesize{\@setfontsize\footnotesize{\rt@pt{9}}{\rt@pt{11}}}
\DeclareRobustCommand\scriptsize{\@setfontsize\scriptsize{\rt@pt{8}}{\rt@pt{9.5}}}
\DeclareRobustCommand\tiny{\@setfontsize\tiny{\rt@pt{6}}{\rt@pt{7}}}
\DeclareRobustCommand\large{\@setfontsize\large{\rt@pt{12}}{\rt@pt{14}}}
\DeclareRobustCommand\Large{\@setfontsize\Large{\rt@pt{14.4}}{\rt@pt{18}}}
\DeclareRobustCommand\LARGE{\@setfontsize\LARGE{\rt@pt{17.28}}{\rt@pt{22}}}
\DeclareRobustCommand\huge{\@setfontsize\huge{\rt@pt{20.74}}{\rt@pt{25}}}
\DeclareRobustCommand\Huge{\@setfontsize\Huge{\rt@pt{24.88}}{\rt@pt{30}}}
% The class's own sizes: each of its styles sets its text with \fontsize{<size>}{1em}, which is
% patched to \rt@fontsize, the size through \rt@pt. A class that no longer has one stops the dump.
\newcommand*\rt@fontsize[1]{\fontsize{\rt@pt{\strip@pt\dimexpr#1\relax}}}
\newcommand*\rt@scalestyle[1]{%
  \expandafter\patchcmd\csname#1\endcsname{\fontsize}{\rt@fontsize}{}%
    {\errmessage{rt-awesome-cv: no \string\fontsize\space in \string\\#1}}}
\@for\rt@style:=%
    headerfirstnamestyle,headerlastnamestyle,headerpositionstyle,headeraddressstyle,%
    headersocialstyle,headerquotestyle,footerstyle,sectionstyle,subsectionstyle,%
    paragraphstyle,entrytitlestyle,entrypositionstyle,entrydatestyle,entrylocationstyle,%
    descriptionstyle,subentrytitlestyle,subentrypositionstyle,subentrydatestyle,%
    subentrylocationstyle,subdescriptionstyle,honortitlestyle,honorpositionstyle,%
    honordatestyle,honorlocationstyle,skilltypestyle,skillsetstyle,lettersectionstyle,%
    recipientaddressstyle,recipienttitlestyle,lettertitlestyle,letterdatestyle,%
    lettertextstyle,letternamestyle,letterenclosurestyle%
  \do{\expandafter\rt@scalestyle\expandafter{\rt@style}}
% The space the class puts before a section's title is \acvSectionTopSkip, 3 mm, and parskip's
% space between paragraphs, half a line, which is also what parts one entry from the next; a
% list's bullets have none between them. The space after a title, \acvSectionContentTopSkip, stays
% the class's 2.5 mm at every density, since the class takes 3 mm back right after it in a
% paragraph section, and 2 mm in the others, and at less than the 2.5 mm a paragraph at the
% smaller sizes would be set over its title. \rt@density is a percentage, 100 until \rtdensity
% sets another, and \rt@gap{<length>} is the length at that share, rounded to a hundredth of a
% point. \rt@itemsep is the space between bullets, which cvitems is patched to take, none at
% normal density and a point less or more for every half of it less or more.
\newcommand*\rt@density{100}
\newcommand*\rt@gap[1]{%
  \the\dimexpr\numexpr\dimexpr#1*\rt@density/100\relax*100/65536\relax\p@/100\relax}
\newcommand*\rt@itemsep{0pt}
\patchcmd{\cvitems}{noitemsep}{noitemsep, itemsep=\rt@itemsep}{}%
  {\errmessage{rt-awesome-cv: no noitemsep in \string\cvitems}}
\newcommand*\rt@skips{%
  \setlength\parskip{\rt@gap{\rt@pt{6.8}}\@plus\rt@gap{\rt@pt{2}}}}
% \rtfontsize{<points>}: every size from <points> in place of 11, then the normal size and the
% space between paragraphs, as size11.clo and parskip set them as the class loads.
\newcommand*\rtfontsize[1]{%
  \edef\rt@size{\number\dimexpr#1\p@\relax}%
  \normalsize
  \rt@skips}
% \rtdensity{<percent>}: the space before a section's title and between paragraphs at that share
% of what the class and the size give them, and the space between bullets.
\newcommand*\rtdensity[1]{%
  \edef\rt@density{\number#1}%
  \edef\acvSectionTopSkip{\rt@gap{3mm}}%
  \edef\rt@itemsep{\the\dimexpr\numexpr\rt@density-100\relax\p@/50\relax}%
  \rt@skips}
\makeatother

% A list section's items, which the editor sets through the class's cvitems in a cvparagraph.
% cvitems takes 4 mm back above its list and below it, space that the row before it leaves in
% \cventry's table; outside a table the list would be set over the section's title, so the port
% gives the 4 mm back, and a list's first line falls where a paragraph's does.
\newenvironment{rtlist}
  {\begin{cvparagraph}\vspace{4mm}\begin{cvitems}}
  {\end{cvitems}\vspace{4mm}\end{cvparagraph}}

% The example's colours. The accent is its red, named rtaccent; the accent control defines it
% again after the format, and sets awesome from it again, since \colorlet takes a colour's value
% when it runs.
\definecolor{rtaccent}{HTML}{DC3522}
\colorlet{awesome}{rtaccent}
\setbool{acvSectionColorHighlight}{true}

% The example's separator between the header's contacts, a bar, which is the class's own.
\renewcommand{\acvHeaderSocialSep}{\quad\textbar\quad}
